\begin{tikzpicture}[every node/.style={font=\small}]

    %==========================================================
    % Bayesian ARX planning model with output memory M_y = 1,
    % unrolled over the horizon tau = (t+1, ..., t+K).
    %   predictive factor : p(y_k | u_k, ybar_k, D_t) = T(M_t^T phi(x_k), ...)
    %   control prior     : p(u_k)
    %   goal prior        : p_*(y_k)
    % The measured output y_t enters as a clamped edge.
    % All branch edges are strictly vertical.
    %==========================================================

    % ---- output row -----------------------------------------
    \node [style=observation]                                       (yt)  at (0.0, 0) {};
    \node [draw, rectangle, minimum width=6mm, minimum height=6mm]  (fy1) at (1.5, 0) {$f_y$};
    \node [style=deterministic]                                     (eq1) at (2.8, 0) {$=$};
    \node [draw, rectangle, minimum width=6mm, minimum height=6mm]  (fy2) at (4.1, 0) {$f_y$};
    \node [style=deterministic]                                     (eq2) at (5.4, 0) {$=$};
    \node [draw, rectangle, minimum width=6mm, minimum height=6mm]  (fy3) at (7.2, 0) {$f_y$};
    \coordinate                                                      (yT)  at (8.5, 0);

    % ---- goal-prior factors terminating the output edges -----
    \node [style=stochastic, minimum width=6mm, minimum height=6mm] (g1) at (2.8, -1.3) {$f_*$};
    \node [style=stochastic, minimum width=6mm, minimum height=6mm] (g2) at (5.4, -1.3) {$f_*$};
    \node [style=stochastic, minimum width=6mm, minimum height=6mm] (g3) at (8.5, -1.3) {$f_*$};

    % ---- control priors, above each predictive factor --------
    \node [style=stochastic, minimum width=6mm, minimum height=6mm] (fu1) at (1.5, 1.3) {$f_u$};
    \node [style=stochastic, minimum width=6mm, minimum height=6mm] (fu2) at (4.1, 1.3) {$f_u$};
    \node [style=stochastic, minimum width=6mm, minimum height=6mm] (fu3) at (7.2, 1.3) {$f_u$};

    %==========================================================
    % Edges
    %==========================================================

    % output chain, with ellipsis between y_{t+2} and the end of the horizon
    \draw[-] (yt)  -- (fy1);
    \draw[-] (fy1) -- (eq1);
    \draw[-] (eq1) -- (fy2);
    \draw[-] (fy2) -- (eq2);
    \draw[-] (eq2) -- (6.0, 0);
    \node at (6.3, 0) {$\cdots$};
    \draw[-] (6.6, 0) -- (fy3);
    \draw[-] (fy3) -- (yT);

    % output edge labels
    \node [style=edgelabel, above] at ($(yt)!0.5!(fy1)$)       {$y_t$};
    \node [style=edgelabel, above] at ($(fy1)!0.5!(eq1)$)      {$y_{t+1}$};
    \node [style=edgelabel, above] at ($(fy2)!0.5!(eq2)$)      {$y_{t+2}$};
    \node [style=edgelabel, above] at ($(fy3.east)!0.5!(yT)$)  {$y_{t+K}$};

    % goal branches
    \draw[-] (eq1) -- (g1);
    \draw[-] (eq2) -- (g2);
    \draw[-] (yT)  -- (g3);

    % control branches into the predictive factors
    \draw[-] (fu1) -- (fy1);
    \draw[-] (fu2) -- (fy2);
    \draw[-] (fu3) -- (fy3);
    \node [style=edgelabel, right=1mm] at ($(fu1.south)!0.5!(fy1.north)$) {$u_{t+1}$};
    \node [style=edgelabel, right=1mm] at ($(fu2.south)!0.5!(fy2.north)$) {$u_{t+2}$};
    \node [style=edgelabel, right=1mm] at ($(fu3.south)!0.5!(fy3.north)$) {$u_{t+K}$};

    %==========================================================
    % Messages: forward predictive messages, the goal message and a
    % backward message that forms the intermediate goal of slice t+1
    %==========================================================
    \msgcircle{down}{right}{fy1.east}{eq1.west}{0.5}{1}
    \msgcircle{down}{right}{fy2.east}{eq2.west}{0.5}{2}
    \msgcircle{right}{up}{g3.north}{yT}{0.5}{3}
    \msgcircle{down}{left}{eq1.east}{fy2.west}{0.5}{4}

\end{tikzpicture}
